\documentclass{article}
\usepackage[T1]{fontenc}
\usepackage{lmodern}
\usepackage{graphicx}
\usepackage{amsmath,amssymb}
\usepackage[a4paper,margin=2cm]{geometry}
\usepackage[absolute,overlay]{textpos}
\setlength{\TPHorizModule}{1mm}
\setlength{\TPVertModule}{1mm}
\usepackage[indonesian]{babel}
\usepackage{hyperref}
\setlength{\emergencystretch}{2em}
% unit: Halaman 20

\begin{document}

% === Halaman 20 (python/native) ===

Selanjutnya untuk setiap putaran 1 sampai r dilakukan operasi 

XOR, pergeseran ke kiri secara sirkuler, dan penjumlahan dalam 

modulo 232 dengan kunci putaran sebagai berikut:


for i  1 to r do


      A  ((A  B) <<< B) + S[2i]


      B  ((B  A) <<< A) + S[2i+1]


endfor

20

\end{document}
